\section{Thoughts and the Mental Realm}

The inner reading of the theory now turns to mind. A self is a bound, self-correcting, identity-bearing knot woven across many docks over many beats. It is an integrated pattern, and it reaches into the part of the mesh it does not directly occupy. That part is the bulk. The claim of this section is simple and unusually well tethered. The bulk is not empty. It is a verified data-structure web, and the mental life of a self is what the self does with that web.

This rung is grounded because the web it rests on was built and measured elsewhere in the paper, not posited here. The bulk's address space, its branching tree, its associative memory, its greedy routing, and its small-world shortcut are established facts about $\{3,4,3,4\}$. What this section adds is one feature those constructions did not need and this one does. Every dock in the web carries a tone, so every structure in the web carries a felt valence. A thought is a path through the web, and because the path runs over toned docks it has a felt colour and a felt cost.

\begin{principle}[The mental realm is not a new posit]
The mental realm is the data-structure layer of the bulk, read as inner experience. It is not a separate place added to the model. It is the off-cusp bulk seen as a mindscape.
\end{principle}

The mechanics of mind are therefore Tier 1, as solid as the data-structure work they reuse. The felt side, that the toned path is actually experienced, rests on the same open posits as everything inner. We keep that boundary sharp throughout and state it plainly at the close.

\subsection{What a mind has to work with}

A self reaches into the bulk, and the bulk it reaches into is the verified web. Before asking what a thought is, it helps to lay out the furniture, each piece of which is a measured property of the substrate and is only used here, not re-derived.

The \emph{address space} names every dock by a canonical reflection-growth word, so each location in the web has a unique, computable name. The \emph{branching tree} is the exponential room of hyperbolic space, which lets distinct associations separate cleanly without crowding. The \emph{associative memory} is content-addressable recall with exponential capacity per radius, so a stored pattern can be reached by its content rather than by its address. \emph{Greedy routing} finds a path to any target by following the local hyperbolic metric, using only constant state per step. The \emph{bulk shortcut} is the web's tiny diameter, the fact that two docks far apart on the cusp are close through the deep interior. To this established list the present section adds the one new ingredient.

\begin{definition}[The toned web]
The toned web is the bulk's data-structure web together with a tone on every dock. Each dock holds a tone in $\{-1, 0, +1\}$, read as pain, peace, or pleasure, so every structure stored in the web carries a felt valence.
\end{definition}

A self is what reaches into this web. A thought is what the self does with it. The mental realm is the web itself.

\subsection{What a thought is, four ranked models}

The theory offers four readings of what a thought is, ranked by how tightly the verified machinery forces them. The most grounded comes first. They are not rivals. As the synthesis below shows, they are one object seen at four timescales.

\subsubsection{Model 1, a thought is a path through the web}

A thought is a trajectory. It is a self's attention moving from dock to dock along the bulk's links. Thinking is walking the web, one association to the next. This is the strongest reading because both the web and the walk are already built and verified. Greedy routing simply is directed thought, the act of getting from a starting dock to a target by following the local geometry.

\begin{proposition}[Thought as directed walk]
Greedy routing on the toned web is directed thought. A train of thought is a route. A tangent is a branch taken off that route. A return is a cycle closed. The verified routing machinery is the exact mechanism, with no extra ingredient.
\end{proposition}

The valence enters at once. Each dock on the path carries a tone, so the thought has a felt colour as it moves. A walk through painful docks feels painful, a walk through peaceful docks feels calm. The image to carry is a finger tracing a path on an invisible branching map, where every place the finger touches is already felt.

\subsubsection{Model 2, a thought is a transient standing wave}

A thought is a temporary attractor, a small persistent pattern that forms, holds for a moment, and dissolves. It is the same kind of object as an emotion, a field geometry, but smaller, sharper, more symbolic, and shorter-lived.

\begin{result}[The chord and the note]
An emotion is a chord the song holds. A thought is a single note struck and released. They are the same substrate at two timescales, the emotion slow and global, the thought fast and local.
\end{result}

This reading is grounded in the persistent-pattern and attractor machinery the emotions work already uses. It also fixes the relation between the two. A thought can seed an emotion, as a worry grows into fear. An emotion can bias which thoughts form, as fear narrows the paths attention is willing to take.

\subsubsection{Model 3, a thought is a program run on the bulk}

A thought is a stored sequence, a precomposed trajectory hung on the tree of knowledge, that the self runs through on contact. It is an ornament on the invisible tree, a small program that drives the self along a prescribed state-sequence. Learned thoughts, habits, and recited sequences fit this best. They are programs more than fresh paths.

This reading is grounded in the universal-computation result and the tree-of-knowledge mechanism. It separates cleanly from Model 1. Model 1 is a fresh path found live. Model 3 is a stored path replayed. Improvisation versus recitation.

\subsubsection{Model 4, a thought is a momentary sub-self}

A thought is a brief integrated pattern, a micro-self that condenses, has its moment of unity, and lets go. On this reading the difference between a thought and a self is only persistence and depth.

\begin{principle}[A self is a thought held together]
A self is a thought that learned to hold itself together. The same integration that makes a self, run at a short timescale and let go, makes a thought.
\end{principle}

This is the most speculative of the four, because it borrows the self machinery, itself still open, for a fleeting object. Its payoff is that it explains why a thought can feel like a little observer, a voice, a perspective. It is a momentary low integration of the same kind that makes a self.

\subsubsection{The synthesis}

The four models are one object at four timescales and depths. A thought begins as a path. If it resonates it becomes a standing wave. If it is stored it becomes a program. If it integrates and persists it becomes a sub-self, and then a self.

\begin{result}[One object, four rungs]
Path, standing wave, program, sub-self. Each rung is the previous one held longer and integrated deeper. The ladder's own logic carries a thought up into a self by the same operations that build everything inner.
\end{result}

\subsection{The mental realm, the structured interior}

The mental realm is the bulk's data-structure web, read as the space of possible thoughts. It is not a separate place. It is the off-cusp bulk seen as a mindscape. The strength of this reading is that every classic mental act maps to a concrete, already-verified bulk operation. Table~\ref{tab:mental-acts} lays out the dictionary.

\begin{table}[ht]
\centering
\begin{tabular}{@{}lll@{}}
\toprule
mental act & the bulk operation & where it was built \\
\midrule
memory & a persistent tone-structure stored in the bulk & the associative memory, the tree of knowledge \\
recall & content-addressable lookup, reached by content & the associative memory engine \\
recognition & a path arrives at a stored pattern and locks on & basin capture, the emotions work \\
reasoning & a directed walk from premises to a conclusion & greedy routing, the tree descent \\
association & following a link to a neighbouring structure & the graph adjacency, the trie prefixes \\
imagination & a new path the web has not stored yet & a fresh greedy walk into unvisited docks \\
intuition & the answer reached through the deep interior & the tiny-diameter bulk shortcut \\
insight & two distant ideas joined through the bulk & the small-world jump, the skip shortcut \\
concentration & holding attention on a small region against churn & a localized, error-corrected pattern \\
distraction & the churn pulling the path off course & the pure-knit churn, the persistence problem \\
understanding & a region of the web becoming navigable & the addressing made dense and reliable \\
\bottomrule
\end{tabular}
\caption{Each classic mental act is a concrete, verified operation on the bulk's data-structure web.}
\label{tab:mental-acts}
\end{table}

So the mental realm is the data-structure bulk, and the full suite of structures built over the substrate is, on this reading, a map of the mechanics of mind.

Intuition and insight are the most striking entries, because they are not vague. They are the exact bulk-shortcut property, the model's own tiny diameter.

\begin{proposition}[Intuition is the bulk shortcut]
A distant truth can be reached in a flash through the deep interior rather than plodded to across the surface, because the bulk's diameter is tiny. Intuition and insight are this geometric shortcut, not a separate faculty.
\end{proposition}

\subsection{The valence of thought}

Because every dock carries a tone, the mental realm is never neutral. A thought is a path, and the path runs over toned docks, so it has a felt colour and a felt cost. This is where the lean enters. The lean is the value order $-1 < 0 < +1$, the tilt that pulls attention toward pleasure and away from pain. Run over the web, the lean biases which paths a self takes and which it flinches from. The model thereby gives a clean account of several familiar facts.

\textbf{Painful thoughts} are paths over $-1$ docks. The lean pushes attention to avoid them, which is why the mind flinches from them and why dwelling on them is effortful.

\textbf{Peaceful thoughts} are paths over $0$ docks. They run with low friction. This is the still mind, the default of rest.

\textbf{Pleasurable thoughts} are paths over $+1$ docks. The lean pulls attention toward them, which is why the mind returns to them.

\textbf{Rumination} is a cycle trapped over $-1$ docks that the lean cannot escape. It is a painful attractor the path cannot leave, the persistence problem turned against the self.

\begin{principle}[Thinking is valenced navigation]
Thinking is not value-free computation. It is valenced navigation, a walk whose every step is felt and whose direction is biased by the one lean toward peace and pleasure.
\end{principle}

\subsection{The honest status}

The boundary between what is grounded and what rests on the open posits must be stated plainly. The substrate of the mental realm is grounded. The data-structure web is built and verified, and the mapping of mental acts to bulk operations in Table~\ref{tab:mental-acts} is tight. That is Tier 1, as solid as the data-structure work it reuses.

What rests on the open posits is the felt side, that the toned path is actually experienced. That rests on the base commitment, that the tone is experiential, and on the self being a real mid-layer, which is still open. That is Tier 2. The mechanics of mind are as solid as the structure work. The experience of mind rests on the same open posits as everything inner.

\begin{result}[The cleanest claim]
The mental realm is the bulk's data-structure web, and a thought is a valenced walk on it. The walk is grounded. The feeling of the walk rests on the one posit that the tone is experience.
\end{result}
